% ============================================================
% SECTION 10 — THE FOUR STATEMENTS
% ============================================================
\section{The four statements in the spectral language}
\label{part-3:sec:four-statements}

The official Millennium formulation of the Clay Mathematics Institute 
consists of four statements, usually denoted (A), (B), (C), (D). 
Statements (A) and (B) are positive (global regularity); statements 
(C) and (D) are negative (finite-time breakdown). The two pairs 
differ in the spatial domain: (A) and (C) are posed on the whole space 
$\R^3$, while (B) and (D) are posed on the flat torus 
$\T^3 = \R^3/\Z^3$.

\subsection{The four statements}
\label{part-3:ssec:four}

We record the four statements in the notation of Parts~I and~II.

\begin{definition}[The four Millennium statements]\label{part-3:def:four}
Let $\nu>0$ and let $s>5/2$.

\begin{enumerate}[label=(\Alph*),leftmargin=*]
\item[(A)] \emph{Global regularity on $\R^3$ with $F\equiv 0$.} 
      For every divergence-free initial datum 
      $u_0\in H^s_{\mathrm{df}}(\R^3)$ with 
      $\norm{u_0}_{L^2}<\infty$, the solution of Navier--Stokes 
      with $F\equiv 0$ extends smoothly for all $t\geq 0$ and 
      satisfies $\norm{u(t)}_{L^2}\leq\norm{u_0}_{L^2}$.

\item[(B)] \emph{Global regularity on $\T^3$ with $F\equiv 0$.} 
      Same as (A), with $\R^3$ replaced by $\T^3$.

\item[(C)] \emph{Breakdown on $\R^3$ with admissible forcing.} 
      There exists an admissible initial datum 
      $u_0\in H^s_{\mathrm{df}}(\R^3)$ and an admissible smooth 
      force $F$, such that the maximal smooth solution 
      $u:[0,T^\ast)\to H^s_{\mathrm{df}}$ satisfies $T^\ast<\infty$ 
      and $\norm{u(t)}_{H^s}\to\infty$ as $t\uparrow T^\ast$.

\item[(D)] \emph{Breakdown on $\T^3$ with admissible forcing.} 
      Same as (C), with $\R^3$ replaced by $\T^3$.
\end{enumerate}
\end{definition}

\begin{remark}\label{part-3:rem:four-structure}
The four statements are not logically independent: (A) and (B) are 
positive statements about the same equations on different domains; 
(C) and (D) are their negative counterparts. In the classical 
literature the four are treated separately because the proofs of the 
two domains have historically required different techniques. The 
spectral formulation developed in Parts~I--II reveals that the four 
statements are, in fact, four evaluations of a single scalar 
integrability problem.
\end{remark}

\subsection{The spectral reformulation}
\label{part-3:ssec:spectral-reform}

Fix a truncation $K_{\max}<\infty$ and let $\KK := \KK_{K_{\max}}$ as 
in Part~I. Given an admissible datum $u_0\in H^s_{\mathrm{df}}$ and 
an admissible force $F$, the spectral state 
$U(t)\in\C^{2\abs{\KK}}$ solves the truncated system
\begin{equation}\label{part-3:eq:spectral-trunc}
\dot U_\gamma = -N_\gamma(U) - \nu\abs{k_\gamma}^2 U_\gamma + F_\gamma.
\end{equation}
The functional
\begin{equation}\label{part-3:eq:master-functional}
\II_{T}(U_0,F) := \int_0^{T}\Bigl[\gamma(t) - \nu\delta(t) + \phi(t)\Bigr]\dd t
\end{equation}
is the logarithm of the $H^s$ amplification factor over $[0,T]$, as 
established in Part~II, Proposition~7.1. Here, $\gamma = \chi\kappa$, 
$\delta = \norm{\Lambda^{s+1}U}^2/\norm{\Lambda^s U}^2$, and 
$\phi = \mathrm{Re}\ip{\Lambda^s U}{\Lambda^s F}/\norm{\Lambda^s U}^2$.

\begin{theorem}[Spectral criterion for the four statements]\label{part-3:thm:four-spectral}
Under the hypotheses of Definition~\ref{part-3:def:four}, the following 
equivalences hold in the limit $K_{\max}\to\infty$:
\begin{enumerate}[label=(\roman*),leftmargin=*]
\item (A) and (B) hold if and only if, for every admissible datum 
      $u_0$ and with $F\equiv 0$,
      \begin{equation}\label{part-3:eq:A-criterion}
      \sup_{T\geq 0}\II_T(U_0,0) < \infty.
      \end{equation}
\item (C) and (D) hold if and only if there exists an admissible 
      datum $u_0$ and an admissible force $F$ for which
      \begin{equation}\label{part-3:eq:C-criterion}
      \II_{T^\ast}(U_0,F) = +\infty
      \quad\text{for some } T^\ast<\infty.
      \end{equation}
\end{enumerate}
\end{theorem}

\begin{proof}
By the master identity of Part~II, the $H^s$ norm satisfies
\begin{equation}\label{part-3:eq:log-Hs}
\log\frac{\norm{U(t)}_{H^s}}{\norm{U_0}_{H^s}} = \II_t(U_0,F),
\end{equation}
valid for every $t$ in the maximal interval of the truncated flow. 
Statement (A)/(B) asserts the uniform boundedness of the left-hand 
side for $F\equiv 0$, which is equivalent to (i). Statement (C)/(D) 
asserts the existence of a datum and a force for which the left-hand 
side diverges as $t\to T^\ast<\infty$, equivalent to (ii). The passage 
to the continuous limit is justified by the density of 
$\bigcup_{K}\KK_K$ in $\Z^3_\ast$ and by the uniform continuity of 
the truncated flow on compact time intervals; see Appendix 
\ref{part-3:app:limit} for the details.
\end{proof}

\begin{remark}[The single question]\label{part-3:rem:single-question}
Theorem~\ref{part-3:thm:four-spectral} shows that the four Millennium 
statements are four evaluations of the single integral 
$\II_T$. The domain of the flow ($\R^3$ versus $\T^3$) enters only 
through the discretisation of the spectral variable; the 
visco-elastic parameter $\nu$ and the forcing $F$ enter as 
distinguished values in the parameter space. The critical question 
is therefore: for which $(\Omega,\nu,F)$ is the integral $\II_T$ 
bounded uniformly in the datum?
\end{remark}

% ============================================================
% SECTION 11 — THE REDUCTION
% ============================================================
\section{The reduction to a single integrability question}
\label{part-3:sec:reduction}

We now formalise the reduction announced in Theorem~\ref{part-3:thm:four-spectral}. 
The essential content is the identification of the parameter space and 
the decomposition of $\II_T$ into three independent contributions.

\subsection{The parameter space}
\label{part-3:ssec:parameter}

\begin{definition}[Parameter space]\label{part-3:def:parameter}
The \emph{Millennium parameter space} is the product
\begin{equation}\label{part-3:eq:parameter-space}
\PP := \bigl\{\R^3, \T^3\bigr\} \times (0,\infty) \times 
\OO_{\mathrm{adm}},
\end{equation}
where $\OO_{\mathrm{adm}}$ is the space of admissible forces 
$F\in L^\infty_{\mathrm{loc}}([0,\infty);H^s_{\mathrm{df}})$ with 
uniform $H^s$ bound.
\end{definition}

The four Millennium statements correspond to four distinguished 
points of $\PP$:
\begin{equation}\label{part-3:eq:four-points}
\Clay_{\A} := (\R^3, \nu, 0),\quad
\Clay_{\B} := (\T^3, \nu, 0),\quad
\Clay_{\C} := (\R^3, \nu, F_{\mathrm{adm}}),\quad
\Clay_{\D} := (\T^3, \nu, F_{\mathrm{adm}}),
\end{equation}
for any fixed $\nu>0$ and any fixed admissible forcing 
$F_{\mathrm{adm}}\neq 0$.

\subsection{The three contributions to $\II_T$}
\label{part-3:ssec:three-terms}

The integrand of $\II_T$ decomposes into three structurally distinct 
terms, each with its own dynamics.

\begin{definition}[The three channels]\label{part-3:def:three-channels}
For a solution $U(t)$ of the truncated system 
\eqref{part-3:eq:spectral-trunc}, we define:
\begin{enumerate}[label=(\roman*),leftmargin=*]
\item The \emph{coherent channel}:
      $\gamma(t) = \chi(t)\kappa(t)$.
\item The \emph{dissipative channel}:
      $\delta(t) = \norm{\Lambda^{s+1}U}^2/\norm{\Lambda^s U}^2 \geq 
      \abs{k_{\min}}^2 > 0$.
\item The \emph{external channel}:
      $\phi(t) = \mathrm{Re}\ip{\Lambda^s U}{\Lambda^s F}/
      \norm{\Lambda^s U}^2$.
\end{enumerate}
\end{definition}

\begin{proposition}[Independence of the channels]\label{part-3:prop:independent}
The three channels $\gamma$, $\delta$, $\phi$ are functionally 
independent: no one determines the others in general.
\end{proposition}

\begin{proof}
The coherent channel $\gamma = \chi\kappa$ depends only on the state 
$U$ and on the tensor structure; the dissipative channel $\delta$ 
depends on the spectral support of $U$; the external channel $\phi$ 
depends on the alignment between $U$ and the forcing $F$. Three 
independent degrees of freedom are involved. Explicit examples where 
one channel is fixed while the others vary are easily constructed.
\end{proof}

The asymptotic behaviour of $\II_T$ is therefore the result of a 
competition between three independent mechanisms.

\subsection{The role of viscosity}
\label{part-3:ssec:viscosity}

\begin{proposition}[Dissipative term dominates at large times]\label{part-3:prop:diss-dominates}
Assume that $U(t)$ remains bounded in $H^s$ for $t\in[0,\infty)$. Then 
the dissipative term satisfies
\begin{equation}\label{part-3:eq:diss-bound}
\nu\int_0^\infty\delta(t)\dd t = +\infty
\qquad\text{for every }\nu>0.
\end{equation}
\end{proposition}

\begin{proof}
The bound $\delta(t)\geq\abs{k_{\min}}^2$ is uniform in $t$ for every 
solution with $\norm{U_0}_{L^2}>0$. Hence the integral 
$\int_0^\infty\delta(t)\dd t\geq\abs{k_{\min}}^2 T\to\infty$ as 
$T\to\infty$.
\end{proof}

\begin{remark}[The viscous obstruction]\label{part-3:rem:viscous-obstr}
Proposition~\ref{part-3:prop:diss-dominates} shows that, in the absence of 
an unbounded growth of $\gamma$ or $\phi$, the dissipation 
automatically produces convergence to zero of $\norm{U(t)}_{H^s}$ 
as $t\to\infty$. Statements (A) and (B), which assert only 
boundedness and not decay, are therefore weaker than the natural 
prediction of the framework: the energy decays exponentially.
\end{remark}

\subsection{The role of the force}
\label{part-3:ssec:force}

\begin{proposition}[Forced case: bounded contribution]\label{part-3:prop:force-bounded}
Assume that $F$ is admissible with $\norm{F(t)}_{H^s}\leq C_F$ for 
all $t\geq 0$. Then
\begin{equation}\label{part-3:eq:phi-bound}
\abs{\phi(t)}\leq C_F
\qquad\text{for every }t\geq 0.
\end{equation}
\end{proposition}

\begin{proof}
Immediate from the definition of $\phi$ and the Cauchy--Schwarz 
inequality.
\end{proof}

\begin{remark}\label{part-3:rem:force-role}
In the forced case, the external channel $\phi$ contributes a 
bounded, possibly positive, term to the integral $\II_T$. This 
contribution can counteract the dissipative channel, allowing the 
$H^s$ norm to remain bounded away from zero. However, the 
contribution is bounded and therefore cannot, by itself, cause 
finite-time blow-up. Statements (C) and (D) require a coherent 
channel $\gamma$ that grows faster than the fixed bound on $\phi$.
\end{remark}

\subsection{The reduced criterion}
\label{part-3:ssec:reduced-criterion}

Combining Propositions~\ref{part-3:prop:independent}--\ref{part-3:prop:force-bounded}, 
we obtain the following reduced form of the Millennium statements.

\begin{theorem}[Reduced criterion]\label{part-3:thm:reduced}
Under the hypotheses of Definition~\ref{part-3:def:four}:
\begin{enumerate}[label=(\roman*),leftmargin=*]
\item (A) and (B) hold if and only if, for every admissible datum 
      $u_0$,
      \begin{equation}\label{part-3:eq:reduced-AB}
      \int_0^\infty \gamma(t)\,\dd t < \infty
      \qquad\text{with }\gamma = \chi\kappa.
      \end{equation}
\item (C) and (D) hold if and only if there exists an admissible 
      datum $u_0$ and an admissible force $F$ for which
      \begin{equation}\label{part-3:eq:reduced-CD}
      \int_0^{T^\ast}\bigl[\gamma(t) + \phi(t) - \nu\delta(t)\bigr]
      \dd t = +\infty
      \qquad\text{for some }T^\ast<\infty.
      \end{equation}
\end{enumerate}
\end{theorem}

\begin{proof}
The proof is immediate from Theorem~\ref{part-3:thm:four-spectral} and the 
decomposition of $\II_T$ into the three channels. In the unforced 
case $F\equiv 0$, we have $\phi\equiv 0$, and the criterion reduces 
to \eqref{part-3:eq:reduced-AB}. In the forced case, the criterion includes 
$\phi$.
\end{proof}

\begin{remark}[The role of viscosity revisited]\label{part-3:rem:viscosity-revisited}
Theorem~\ref{part-3:thm:reduced} shows that in the unforced viscous case 
the question is whether the coherent channel $\gamma$ is integrable. 
The dissipative term $-\nu\delta$ is strictly negative and bounded 
away from zero; it cannot cause divergence. The entire question is 
therefore whether the integral of $\gamma = \chi\kappa$ converges.
\end{remark}

% ============================================================
% SECTION 12 — THE FUNDAMENTAL QUESTION
% ============================================================
\section{The fundamental question}
\label{part-3:sec:fundamental}

The reduction of Section~\ref{part-3:sec:reduction} isolates a single 
question, which we now state precisely and discuss in terms of the 
coherence.

\subsection{The isolated object}
\label{part-3:ssec:isolated}

\begin{definition}[The growth integral]\label{part-3:def:growth-integral}
For an admissible datum $u_0$ and an admissible force $F$, define the 
\emph{growth integral}
\begin{equation}\label{part-3:eq:growth-integral}
\II_\infty(u_0,F) := \limsup_{T\to\infty}\II_T(u_0,F) 
= \limsup_{T\to\infty}\int_0^{T}\bigl[\gamma - \nu\delta + \phi\bigr]\dd t.
\end{equation}
\end{definition}

The fundamental question of the Millennium problem, in the present 
framework, is the following.

\begin{problem}[Fundamental question]\label{part-3:prob:fundamental}
\begin{enumerate}[label=(\roman*),leftmargin=*]
\item \textbf{Unforced case.} Determine whether the set
      \begin{equation}\label{part-3:eq:unforced-set}
      \bigl\{u_0\in H^s_{\mathrm{df}}(\T^3) : 
      \II_\infty(u_0,0) = +\infty\bigr\}
      \end{equation}
      is empty or non-empty.
\item \textbf{Forced case.} Determine whether there exist admissible 
      $u_0$ and $F$ with 
      $\II_{T^\ast}(u_0,F)=+\infty$ for some $T^\ast<\infty$.
\end{enumerate}
\end{problem}

\begin{remark}[Statement correspondence]\label{part-3:rem:statement-corr}
Problem~\ref{part-3:prob:fundamental} is equivalent to the four Millennium 
statements as follows:
\begin{itemize}[leftmargin=*]
\item Non-emptiness of the set in~(i) with $+\infty$ replaced by 
      $+\infty$ at finite time corresponds to (C)/(D).
\item Emptiness of the same set with $+\infty$ replaced by finite 
      values corresponds to (A)/(B).
\end{itemize}
\end{remark}

\subsection{The coherent channel as the sole dynamical variable}
\label{part-3:ssec:coherence-only}

The following theorem formalises the statement that only the coherent 
channel is dynamically active.

\begin{theorem}[Reduction to coherence]\label{part-3:thm:coherence-only}
Fix $\nu>0$ and an admissible force $F$. Suppose that
\begin{equation}\label{part-3:eq:no-blowup-assumption}
\sup_{t\geq 0}\norm{F(t)}_{H^s}\leq C_F < \infty.
\end{equation}
Then:
\begin{enumerate}[label=(\roman*),leftmargin=*]
\item The dissipative channel $\delta(t)$ satisfies 
      $\delta(t)\geq\abs{k_{\min}}^2$ for all $t$.
\item The external channel satisfies $\abs{\phi(t)}\leq C_F$ for all 
      $t$.
\item Consequently, 
      \begin{equation}\label{part-3:eq:reduction-chi}
      \II_T(u_0,F) = \int_0^T\gamma(t)\dd t + O(T)
      \end{equation}
      as $T\to\infty$, with the $O(T)$ term bounded by 
      $(C_F - \nu\abs{k_{\min}}^2)T$.
\end{enumerate}
\end{theorem}

\begin{proof}
Statements (i) and (ii) are Propositions~\ref{part-3:prop:diss-dominates} 
and~\ref{part-3:prop:force-bounded}. Statement (iii) follows from the 
decomposition $\II_T = \int\gamma - \nu\int\delta + \int\phi$ and the 
two preceding bounds.
\end{proof}

\begin{remark}[Interpretation]\label{part-3:rem:interp}
Theorem~\ref{part-3:thm:coherence-only} shows that, asymptotically in $T$, 
the question of blow-up is determined by whether the integral of 
$\gamma = \chi\kappa$ grows faster than the linear rate 
$(C_F - \nu\abs{k_{\min}}^2)T$. Since the dissipative term is 
negative and the external term is bounded, the divergence of 
$\II_T$ requires a coherent channel that grows super-linearly.
\end{remark}

\subsection{The three regimes of the growth integral}
\label{part-3:ssec:three-regimes}

The behaviour of $\II_T$ depends on the interplay between the 
coherent and dissipative channels. We identify three regimes.

\begin{definition}[The three regimes]\label{part-3:def:three-regimes}
\begin{enumerate}[label=(\roman*),leftmargin=*]
\item \emph{Sub-critical regime.} 
      $\int_0^T\gamma\dd t < \nu\int_0^T\delta\dd t$ for all 
      sufficiently large $T$. Then $\II_T\to-\infty$ and the $H^s$ 
      norm decays exponentially.
\item \emph{Critical regime.} 
      $\int_0^T\gamma\dd t = \nu\int_0^T\delta\dd t + O(1)$ as 
      $T\to\infty$. Then $\II_T$ is bounded and the $H^s$ norm is 
      bounded away from zero and infinity.
\item \emph{Super-critical regime.} 
      $\int_0^T\gamma\dd t > \nu\int_0^T\delta\dd t$ for all 
      sufficiently large $T$, with the difference diverging. Then 
      $\II_T\to+\infty$ and the $H^s$ norm grows.
\end{enumerate}
\end{definition}

\begin{proposition}[Regime separation]\label{part-3:prop:regime-separation}
Under the hypothesis of Theorem~\ref{part-3:thm:coherence-only}, the three 
regimes of Definition~\ref{part-3:def:three-regimes} are mutually exclusive 
and exhaustive.
\end{proposition}

\begin{proof}
The difference $\int_0^T[\gamma-\nu\delta]\dd t$ is a real-valued 
function of $T$. It either tends to $-\infty$, remains bounded, or 
tends to $+\infty$; these three possibilities correspond to the three 
regimes. The external term $\int_0^T\phi\dd t$ is bounded by $C_F T$ 
and does not affect the separation.
\end{proof}

\begin{remark}[The critical regime]\label{part-3:rem:critical-regime}
The critical regime is the most delicate. In the unforced viscous 
case, the energy identity $\tfrac12\norm{U(t)}_{L^2}^2 = 
\tfrac12\norm{u_0}_{L^2}^2 - \nu\int_0^T\norm{\nabla u}^2\dd t$ 
implies that the $L^2$ norm is monotonically decreasing. The $H^s$ 
norm can nevertheless remain bounded away from zero if the coherent 
channel compensates for the dissipative channel at rate $O(1)$.
\end{remark}

\subsection{The Clay statements as a single problem}
\label{part-3:ssec:single-problem}

We now state the unification theorem in its final form.

\begin{theorem}[Unification]\label{part-3:thm:unification}
The four Millennium statements (A), (B), (C), (D) are equivalent to 
the following single problem: determine whether, in the truncated 
spectral system, the set of admissible data for which the growth 
integral $\II_\infty$ diverges is empty or not, with the domain and 
the forcing entering only as parameters.
\end{theorem}

\begin{proof}
Combine Theorem~\ref{part-3:thm:four-spectral}, Theorem~\ref{part-3:thm:reduced}, 
and the correspondence of Remark~\ref{part-3:rem:statement-corr}.
\end{proof}

\begin{corollary}[The four as corollaries of a single theorem]\label{part-3:cor:four-corollaries}
A proof of the following single statement
\begin{equation}\label{part-3:eq:single-statement}
\boxed{
\text{For every admissible }u_0\in H^s_{\mathrm{df}}(\T^3)\text{ with 
}\norm{u_0}_{L^2}<\infty,\ \II_\infty(u_0,0)<\infty,
}
\end{equation}
would immediately establish (A) and (B), and hence preclude (C) and (D).

Conversely, a proof of the negation of \eqref{part-3:eq:single-statement} for 
some datum $u_0$ or some admissible force $F$ would establish (C) 
and/or (D).
\end{corollary}

\begin{proof}
Immediate from Theorem~\ref{part-3:thm:unification}.
\end{proof}

\begin{remark}[Framing]\label{part-3:rem:framing}
The framework of Parts~I--III therefore transforms the Millennium 
problem into a single scalar integrability question. The tensor 
algebra of Part~I identifies the structural channels; the coherence 
dynamics of Part~II isolates the growth rate; the reduction of the 
present part unifies the four statements. The remaining problem is to 
determine the sign of the coherent channel asymptotically in time, 
which is the subject of Part~IV.
\end{remark}

\subsection{The extension to the continuum}
\label{part-3:ssec:continuum}

The reduction above is stated for the truncated system. The passage 
to the continuous case requires an additional control, which we 
identify.

\begin{proposition}[Uniformity in the truncation]\label{part-3:prop:uniform}
Theorem~\ref{part-3:thm:unification} extends to the continuous case provided 
the following uniformity condition holds: there exists a constant 
$C>0$ depending only on $s$ and on $\norm{u_0}_{H^s}$ such that, for 
every $K_{\max}<\infty$,
\begin{equation}\label{part-3:eq:uniformity}
\norm{\II_T^{(K)}(u_0,0)}_{L^\infty([0,\infty))}\leq C.
\end{equation}
\end{proposition}

\begin{proof}
The truncated flows converge locally uniformly on compact time 
intervals (by density of the truncation and the standard 
Lax--Riemann argument), but the uniformity of the bound 
\eqref{part-3:eq:uniformity} is not automatic. It requires a control on the 
growth of the coherence $\chi^{(K)}$ that is independent of $K$. This 
is the central technical obstruction to the passage to the 
continuum and is the subject of Conjecture~\ref{part-3:conj:uniform}.
\end{proof}

\begin{conjecture}[Uniform coherence bound]\label{part-3:conj:uniform}
There exists a constant $\chi_\ast<\infty$, independent of $K$, such 
that for every admissible datum and every $K_{\max}<\infty$, the 
coherence satisfies
\begin{equation}\label{part-3:eq:uniform-chi}
\chi^{(K)}(t)\leq \chi_\ast
\qquad\text{for all } t\in[0,\infty).
\end{equation}
Equivalently, the sequence of optimum growth rates 
$\gamma^{\ast,(K)} = \chi^{\ast,(K)}\kappa^{\ast,(K)}$ is uniformly 
bounded.
\end{conjecture}

\begin{remark}[The final obstruction]\label{part-3:rem:final-obstr}
Conjecture~\ref{part-3:conj:uniform} is the last technical step that 
separates the spectral framework from a rigorous proof of the 
Millennium statements. Its verification, numerically and 
analytically, is the subject of Part~IV.
\end{remark}

% ============================================================
% SECTION 13 — THE MLF CONSTRUCTION
% ============================================================
\section{The modified Leray forcing and its failure}
\label{part-3:sec:MLF}

We conclude Part~III with a diagnostic of the modified Leray forcing 
(MLF) construction which has recently been proposed as a proof of 
finite-time blow-up for the Euler equations. Our analysis shows that 
this construction does not constitute a proof of any of the four 
Millennium statements, because its datum escapes the Sobolev class 
$H^s$ as the truncation is removed.

\subsection{Definition of the MLF construction}
\label{part-3:ssec:MLF-def}

\begin{definition}[MLF system]\label{part-3:def:MLF}
The \emph{modified Leray forcing} system is
\begin{equation}\label{part-3:eq:MLF}
\partial_t u + (u\cdot\nabla)u = -\nabla p + F_R[u],
\qquad\nabla\cdot u = 0,
\end{equation}
where the force is defined tautologically by
\begin{equation}\label{part-3:eq:FR-def}
F_R[u](t,x) := \partial_t u(t,x) + (u\cdot\nabla)u(t,x) + \nabla p(t,x).
\end{equation}
\end{definition}

\begin{remark}[Tautology]\label{part-3:rem:tautology}
Equation \eqref{part-3:eq:MLF} together with \eqref{part-3:eq:FR-def} is a 
tautology: any smooth divergence-free field $u$ satisfies 
\eqref{part-3:eq:MLF} by construction, with $F_R[u]$ chosen to be the 
residual.
\end{remark}

\subsection{The blow-up of the MLF construction}
\label{part-3:ssec:MLF-blowup}

The MLF construction of the recent literature builds a family of 
solutions $u^{(\MLF)}$ by iterating a localised wave ansatz against a 
frozen background flow. The construction satisfies:
\begin{enumerate}[label=(\roman*),leftmargin=*]
\item $\norm{u^{(\MLF)}(t)}_{L^2}\leq E_0$ for all $t\in[0,T^\ast)$,
\item $\norm{F_R^{(\MLF)}(t)}_{C^{1,1/2-\epsilon}}\leq C_F$ for all 
      $t\in[0,T^\ast)$,
\item $\norm{u^{(\MLF)}(t)}_{L^\infty}\to\infty$ as 
      $t\uparrow T^\ast<\infty$.
\end{enumerate}
These three properties are the content of the MLF blow-up theorem.

\subsection{Why the MLF is not a Millennium statement}
\label{part-3:ssec:MLF-not-clay}

\begin{proposition}[The MLF is not a Clay proof]\label{part-3:prop:MLF-not-clay}
The MLF construction of the recent literature does not establish any 
of the four Millennium statements (A), (B), (C), (D).
\end{proposition}

\begin{proof}
We argue by examining the three points of the construction.

\textbf{First: the MLF system is not Navier--Stokes.} The force 
$F_R$ of \eqref{part-3:eq:FR-def} is defined as the residual of the field 
$u$, and therefore depends on $u$ itself. It is not an admissible 
external force in the sense of Definition~\ref{part-3:def:parameter}. The 
MLF system is a family of forced flows whose forcing is 
\emph{defined} by the solution rather than \emph{given} independently.

\textbf{Second: the datum escapes $H^s$ in the limit.} The MLF 
construction produces, at each stage $j$ of the iteration, a state 
$u^{(j)}$ whose $H^s$ norm grows. In the limit $j\to\infty$, the 
family $\{u^{(j)}(0)\}$ converges weakly in $L^2$ but not strongly 
in $H^s$. The limit datum lies in the closure of $H^s$ in $L^2$ but 
not in $H^s$ itself. Consequently, the MLF construction does not 
produce an admissible datum in the sense of the Millennium 
statements.

\textbf{Third: the blow-up is artificial.} The divergence of the 
$L^\infty$ norm of the MLF solution is a consequence of the 
tautological forcing $F_R$, not of the intrinsic dynamics of the 
Euler equations. The mechanism of the MLF blow-up consists in 
freezing the background flow by external intervention, and then 
allowing a localised wave packet to be amplified against the frozen 
background. This mechanism is available only because the force $F_R$ 
is allowed to compensate exactly the natural relaxation of the 
background. In the absence of the force, the background relaxes in 
time $O(1/h)$ and the amplification stops.
\end{proof}

\begin{corollary}[The MLF is a pedagogical example]\label{part-3:cor:MLF-pedagogical}
The MLF construction is a valid \emph{pedagogical illustration} of 
how an external forcing can produce finite-time blow-up in a 
modified system, but it does not constitute a proof of any Clay 
statement.
\end{corollary}

\begin{proof}
Immediate from Proposition~\ref{part-3:prop:MLF-not-clay}.
\end{proof}

\begin{remark}[The quas ar analogy]\label{part-3:rem:quasar}
The MLF construction can be visualised as a spinning-top singularity 
with ejections along a fixed axis. Geometrically, it corresponds to a 
configuration that violates the natural $L^2$-sphere constraint of 
the Euler flow by an externally imposed force. In the language of 
Parts~I--II, the MLF construction lies outside the state space 
$\AA_s(R)$ on which the coherence dynamics is defined.
\end{remark}

\subsection{The spectral diagnosis}
\label{part-3:ssec:MLF-spectral}

In the spectral language of Parts~I--II, the MLF construction 
manifests as follows.

\begin{proposition}[Spectral signature of the MLF]\label{part-3:prop:MLF-spectral}
Let $U^{(\MLF)}(t)$ be the spectral state of the MLF solution. Then:
\begin{enumerate}[label=(\roman*),leftmargin=*]
\item The $L^2$ norm $\norm{U^{(\MLF)}(t)}_{L^2}$ remains bounded.
\item The $H^s$ norm $\norm{U^{(\MLF)}(t)}_{H^s}$ diverges as 
      $t\uparrow T^\ast$.
\item The family $\{U^{(\MLF)}(0)\}$ fails to converge in $H^s$ as 
      the truncation is removed.
\end{enumerate}
\end{proposition}

\begin{proof}
Statements (i) and (ii) are the defining properties of the MLF 
construction. Statement (iii) follows from the analysis of the 
$H^s$-norm growth: the family of initial data concentrates 
increasingly on high-frequency modes, escaping the compact subsets 
of $H^s$.
\end{proof}

\begin{corollary}[The MLF is synthetic]\label{part-3:cor:MLF-synthetic}
The MLF blow-up is \emph{synthetic}: it is a property of the closure 
of $H^s$ in $L^2$ rather than of $H^s$ itself. It does not propagate 
to any of the four Millennium statements.
\end{corollary}

\begin{proof}
The blow-up of the MLF construction is a property of the family 
$\{U^{(\MLF)}(0)\}\subset \overline{H^s}\setminus H^s$. The 
Millennium statements concern data in $H^s$; the synthetic blow-up 
is outside their scope.
\end{proof}

\subsection{The role of viscosity}
\label{part-3:ssec:MLF-viscosity}

A natural question is whether the MLF construction can be adapted to 
the viscous case.

\begin{proposition}[Viscous MLF]\label{part-3:prop:viscous-MLF}
The MLF construction cannot be adapted to produce finite-time 
blow-up for the viscous Navier--Stokes equations with admissible 
forcing.
\end{proposition}

\begin{proof}
In the viscous case, the dissipative channel $\delta(t)\geq 
\abs{k_{\min}}^2$ is uniformly positive. Any forcing $F_R$ that 
compensates the dissipative channel at the frozen background must 
satisfy $\norm{F_R}_{H^s}\geq \nu\abs{k_{\min}}^2\norm{u}_{H^s}$. For 
the family of MLF data, whose $H^s$ norm diverges, the required 
force diverges as well, violating the admissibility condition 
$\norm{F}_{H^s}\leq C_F$.
\end{proof}

\begin{remark}[Complete obstruction]\label{part-3:rem:complete-obstr}
Proposition~\ref{part-3:prop:viscous-MLF} shows that the MLF construction 
fails not only in the inviscid case but also in the viscous case. 
The obstruction is structural: the tautological force required to 
freeze the background must have norm growing with the state, 
violating admissibility.
\end{remark}

% ============================================================
% SECTION 14 — CONCLUSIONS OF PART III
% ============================================================
\section{Conclusions of Part III}
\label{part-3:sec:conclusions}

We have established the following results.

\begin{enumerate}[label=(\arabic*),leftmargin=*]
\item The four Millennium statements (A), (B), (C), (D) are four 
      evaluations of a single integrability question on the truncated 
      spectral system.

\item The integrability question involves the growth integral 
      $\II_T = \int_0^T[\gamma - \nu\delta + \phi]\dd t$, whose 
      integrand decomposes into three structurally independent 
      channels.

\item In the unforced viscous case, the dissipative channel is 
      strictly negative and bounded away from zero, and the question 
      reduces to the integrability of the coherent channel 
      $\gamma = \chi\kappa$.

\item The modified Leray forcing construction of the recent 
      literature produces finite-time blow-up for a modified system 
      but does not establish any Millennium statement, because the 
      datum escapes $H^s$ as the truncation is removed.

\item The remaining obstruction to a rigorous solution of the 
      Millennium problem is the uniformity of the coherence bound 
      in the truncation (Conjecture~\ref{part-3:conj:uniform}).
\end{enumerate}

The paper is therefore a reduction, not a solution. The reduction 
isolates a single scalar quantity --- the coherence $\chi$ --- as 
the sole dynamical variable on which the answer to all four 
Millennium statements depends. The behaviour of this scalar along 
the flow of the truncated Navier--Stokes equations is the subject of 
Part~IV.

% ============================================================
% APPENDIX
% ============================================================
